\documentclass[12pt]{article}
\usepackage[utf8]{inputenc}
\usepackage[T1]{fontenc}
\usepackage[margin=1in]{geometry}
\usepackage{setspace}
\usepackage{xcolor}
\usepackage[colorlinks=true, urlcolor=blue!60!black, linkcolor=blue!60!black]{hyperref}
\setlength{\parindent}{2em}
\setlength{\parskip}{1.25em}
\renewcommand{\baselinestretch}{1.0}
\begin{document}
\begin{center}
{\LARGE \textit{Arizona v.\ United States}, 567 U.S.\ 387 (2012)\par}
\vspace{0.8em}
{\large Case Brief\par}
\end{center}
\vspace{1em}

\section{Memory Jogger}

In \href{https://supreme.justia.com/cases/federal/us/567/387/}{\textit{Arizona v.\ United States}}, the Supreme Court struck down three of four challenged provisions of Arizona's S.B.\ 1070 as preempted by federal immigration law---barring state crimes for alien registration failures and unauthorized work and state warrantless arrests for removability---while declining to facially enjoin the ``show me your papers'' status-verification mandate before state courts could construe it.

\section{Detailed Case Facts}

On April 23, 2010, Arizona Governor Jan Brewer signed into law the \href{https://www.law.cornell.edu/wex/support_our_law_enforcement_and_safe_neighborhoods_act_of_2010}{Support Our Law Enforcement and Safe Neighborhoods Act}, widely known as S.B.\ 1070 (and amended shortly thereafter by H.B.\ 2162). Facing acute border-state fiscal, public-safety, and law-enforcement pressures---the record noted hundreds of thousands of annual apprehensions in Arizona, an unauthorized population estimated at nearly six percent of the State, and border-region smuggling warnings along an interstate highway south of Phoenix---the Arizona Legislature enacted S.B.\ 1070 to ``discourage and deter the unlawful entry and presence of aliens and economic activity by persons unlawfully present in the United States,'' establishing an official state policy of ``attrition through enforcement.''

Before the statute's scheduled July 29, 2010 effective date, the United States Department of Justice filed a pre-enforcement facial challenge under the Supremacy Clause, U.S.\ Const.\ art.\ VI, cl.\ 2, seeking to enjoin four core provisions of S.B.\ 1070 as preempted by federal immigration statutes, principally the Immigration and Nationality Act (INA) and the Immigration Reform and Control Act of 1986 (IRCA):

\begin{itemize}
\item \textbf{Section 3} (Ariz.\ Rev.\ Stat.\ Ann.\ \S~13--1509): Created a new state misdemeanor for the ``willful failure to complete or carry an alien registration document'' in violation of federal registration statutes, 8 U.S.C.\ \S\S~1304(e) and 1306(a), while ruling out probation and pardon as sentencing options.
\item \textbf{Section 5(C)} (Ariz.\ Rev.\ Stat.\ Ann.\ \S~13--2928(C)): Made it a state misdemeanor---punishable by up to a \$2,500 fine and six months' imprisonment---for an unauthorized alien to ``knowingly apply for work, solicit work in a public place or perform work as an employee or independent contractor'' in Arizona. Under IRCA, 8 U.S.C.\ \S~1324a, Congress imposed civil and criminal penalties on \textit{employers} who knowingly hire unauthorized workers and expressly preempted state sanctions on employers (other than licensing laws, \S~1324a(h)(2)), but imposed only civil consequences (such as removal and ineligibility for adjustment of status) on unauthorized \textit{employees}, apart from federal fraud crimes under 18 U.S.C.\ \S~1546(b).
\item \textbf{Section 6} (Ariz.\ Rev.\ Stat.\ Ann.\ \S~13--3883(A)(5)): Authorized state and local police officers to arrest a person without a warrant if the officer has probable cause to believe the person ``has committed any public offense that makes the person removable from the United States.'' By contrast, federal law treats unlawful presence as a civil rather than criminal matter, limits warrantless arrests by trained federal immigration officers to situations where the alien ``is likely to escape before a warrant can be obtained'' (8 U.S.C.\ \S~1357(a)(2)), and restricts state performance of immigration-officer functions to formal federal agreements (\S~1357(g)(1)) or federal-state cooperation (\S~1357(g)(10)(B)).
\item \textbf{Section 2(B)} (Ariz.\ Rev.\ Stat.\ Ann.\ \S~11--1051(B)): Required state and local officers conducting a lawful stop, detention, or arrest for another offense to make a ``reasonable attempt,'' when practicable, to determine the person's immigration status with Immigration and Customs Enforcement (ICE) if reasonable suspicion exists that the person is unlawfully present, and mandated that any arrestee have his or her immigration status verified before release. Section 2(B) included three statutory limits: a presumption of lawful presence upon presentation of a valid Arizona driver's license or similar identification; a bar on considering race, color, or national origin except as permitted by the U.S.\ and Arizona Constitutions; and a directive to implement the law consistently with federal immigration and civil-rights laws.
\end{itemize}

The constitutional backdrop rested on the Federal Government's broad authority over immigration derived from Congress's power to ``establish an uniform Rule of Naturalization'' (U.S.\ Const.\ art.\ I, \S~8, cl.\ 4), its inherent sovereign power over foreign relations, and the Supremacy Clause (art.\ VI, cl.\ 2), as elaborated in \href{https://www.law.cornell.edu/constitution-conan/article-1/section-8/clause-1/immigration-related-state-laws}{federal immigration-preemption doctrine}. Notably, the United States challenged S.B.\ 1070 solely on preemption grounds rather than on Fourth Amendment or Equal Protection racial-profiling claims.

\section{Procedural History}

On July 28, 2010, Judge Susan Bolton of the U.S.\ District Court for the District of Arizona granted the United States' motion for a preliminary injunction blocking enforcement of Sections 2(B), 3, 5(C), and 6, while allowing other unchallenged or severable portions of S.B.\ 1070 to take effect (\textit{United States v.\ Arizona}, 703 F.\ Supp.\ 2d 980 (D.\ Ariz.\ 2010)).

Arizona filed an interlocutory appeal to the U.S.\ Court of Appeals for the Ninth Circuit. On April 11, 2011, a divided three-judge panel affirmed the preliminary injunction in full (\textit{United States v.\ Arizona}, 641 F.\ 3d 339 (9th Cir.\ 2011)). Writing for the panel, Judge Richard Paez (joined by Judge John T.\ Noonan, Jr., who also concurred separately) concluded that the United States was likely to succeed on its preemption claims against all four sections. Judge Carlos Bea concurred as to Sections 3 and 5(C) but dissented as to Sections 2(B) and 6.

Rather than seeking en banc review, Arizona petitioned directly for a writ of certiorari. The Supreme Court granted certiorari on December 12, 2011 (565 U.S.\ 1092), heard oral argument on April 25, 2012 (with Paul D.\ Clement arguing for Arizona and Solicitor General Donald B.\ Verrilli, Jr.\ arguing for the United States), and issued its decision on June 25, 2012. As reported on \href{https://www.scotusblog.com/cases/arizona-v-united-states/}{SCOTUSblog} and \href{https://www.oyez.org/cases/2011/11-182}{Oyez}, the Court affirmed in part, reversed in part, and remanded, invalidating Sections 3, 5(C), and 6 while lifting the facial injunction against Section 2(B).

\section{Judicial Votes}

The case was decided by a \textbf{5--3} vote, with \textbf{Justice Elena Kagan recused} (taking no part in the consideration or decision of the case due to her prior service as Solicitor General).

\begin{itemize}
\item \textbf{Majority Opinion (5 Justices):} Justice Anthony M.\ Kennedy delivered the opinion of the Court, joined by Chief Justice John G.\ Roberts, Jr., and Justices Ruth Bader Ginsburg, Stephen G.\ Breyer, and Sonia Sotomayor. (Struck down \S\S~3, 5(C), and 6; reversed the preliminary injunction against \S~2(B).)
\item \textbf{Concurring in Part and Dissenting in Part (Justice Scalia):} Justice Antonin Scalia concurred in upholding \S~2(B) and dissented from the invalidation of \S\S~3, 5(C), and 6 (would have upheld all four provisions).
\item \textbf{Concurring in Part and Dissenting in Part (Justice Thomas):} Justice Clarence Thomas concurred in upholding \S~2(B) and dissented from the invalidation of \S\S~3, 5(C), and 6 (would have upheld all four provisions on textualist grounds).
\item \textbf{Concurring in Part and Dissenting in Part (Justice Alito):} Justice Samuel A.\ Alito, Jr., concurred in upholding \S~2(B) and in invalidating \S~3 under \textit{Hines v.\ Davidowitz}, but dissented from the invalidation of \S\S~5(C) and 6.
\end{itemize}

\section{Holding}

\textbf{Judgment:} Affirmed in part (upholding the preliminary injunction against \S\S~3, 5(C), and 6), reversed in part (vacating the facial preliminary injunction against \S~2(B)), and remanded for further proceedings.

\textbf{Rule Announced:} Under the Supremacy Clause and implied \href{https://www.law.cornell.edu/wex/preemption}{preemption} doctrine, federal immigration law displaces state statutes that (1) intrude into a field occupied exclusively by Congress---such as alien registration---even when the state law mirrors or complements federal standards; or (2) stand as an obstacle to Congress's deliberate statutory balance and federal enforcement discretion, such as criminalizing unauthorized employment where IRCA chose only civil sanctions on employees or granting state officers unilateral warrantless arrest power for civil removability. However, a state mandate requiring officers to verify immigration status with ICE during lawful stops or arrests is not facially preempted prior to state-court construction, provided it does not prolong detention solely to check status.

\textbf{Holding vs.\ Dicta:} The Court's facial invalidation of \S\S~3, 5(C), and 6 and its refusal to facially enjoin \S~2(B) on a pre-enforcement record are binding holdings. By contrast, Justice Kennedy's extended discussion of hypothetical Fourth Amendment and preemption violations if \S~2(B) were applied to prolong stops or delay release (citing \textit{Illinois v.\ Caballes}, 543 U.S.\ 405 (2005), and \textit{Arizona v.\ Johnson}, 555 U.S.\ 323 (2009)), his explicit reservation of whether state officers possess authority to prolong stops on reasonable suspicion of federal immigration \textit{crimes}, and his closing civic reflections on naturalization ceremonies and national immigration discourse are influential dicta and express reservations for future as-applied litigation.

\section{Analysis of Opinions}

\subsection{Majority Opinion (Justice Kennedy)}

Writing for five Justices in the \href{https://www.law.cornell.edu/supremecourt/text/11-182}{opinion of the Court}, Justice Kennedy anchored federal immigration supremacy in both the Naturalization Clause (art.\ I, \S~8, cl.\ 4) and the Federal Government's inherent sovereign power over foreign relations, citing \href{https://supreme.justia.com/cases/federal/us/458/1/}{\textit{Toll v.\ Moreno}}, 458 U.S.\ 1, 10 (1982), and \textit{Hines v.\ Davidowitz}, 312 U.S.\ 52, 64 (1941). Foreign nations must be able to communicate with one national sovereign rather than fifty States regarding the treatment of their nationals, and federal removal enforcement is a civil system characterized by broad executive discretion reflecting humanitarian equities and diplomatic considerations.

After acknowledging Arizona's serious border-related challenges and reciting the classic preemption framework---express preemption, field preemption (\textit{Rice v.\ Santa Fe Elevator Corp.}, 331 U.S.\ 218, 230 (1947)), and conflict/obstacle preemption (\textit{Hines}, 312 U.S.\ at 67; \textit{Crosby v.\ National Foreign Trade Council}, 530 U.S.\ 363, 372--373 (2000))---Kennedy evaluated each provision in turn:

\begin{itemize}
\item \textbf{Section 3 (Field Preemption):} Relying squarely on \textit{Hines}, Kennedy held that Congress has occupied the entire field of alien registration through a ``single integrated and all-embracing system'' (8 U.S.C.\ \S\S~1301--1306). Where Congress occupies a field, even parallel or complementary state penalties are barred. Moreover, \S~3 would strip federal officials of prosecutorial control and create sentencing conflicts by eliminating federal options like probation and pardon.
\item \textbf{Section 5(C) (Obstacle Preemption):} Distinguishing \textit{De Canas v.\ Bica}, 424 U.S.\ 351 (1976)---decided when federal law expressed only a ``peripheral concern'' with unauthorized employment---Kennedy emphasized that IRCA established a comprehensive framework targeting \textit{employers} with civil and criminal sanctions while imposing only civil remedies on unauthorized \textit{employees}. Drawing on IRCA's text, structure, and legislative history (including the 1981 Select Commission report), Kennedy concluded Congress made a deliberate choice not to criminalize unauthorized work by employees. Under \textit{Geier v.\ American Honda Motor Co.}, 529 U.S.\ 861 (2000), IRCA's express preemption clause covering employers (\S~1324a(h)(2)) did not foreclose ordinary obstacle preemption; a ``conflict in technique'' can disrupt Congress's calibrated scheme as much as a conflict in overt policy.
\item \textbf{Section 6 (Obstacle Preemption):} Unlawful presence is generally a civil matter, not a crime (\textit{INS v.\ Lopez-Mendoza}, 468 U.S.\ 1032, 1038 (1984)). By empowering state officers to make unilateral warrantless arrests based on probable cause of a removable offense, \S~6 gave state police broader arrest authority than Congress granted trained federal ICE officers (who need a warrant or likelihood of escape under \S~1357(a)(2)). Federal law allows state officers to perform immigration functions only under formal supervised \S~1357(g) agreements or genuine ``cooperation'' (\S~1357(g)(10)(B)), which does not encompass unilateral state arrests divorced from federal instruction or removal discretion.
\item \textbf{Section 2(B) (No Facial Preemption):} Congress actively encouraged federal-state information sharing through 8 U.S.C.\ \S\S~1357(g)(10)(A), 1373(c), and 1644, and ICE's Law Enforcement Support Center operates 24/7 to answer status inquiries (analogous to Arizona's E-Verify mandate upheld in \href{https://supreme.justia.com/cases/federal/us/563/582/}{\textit{Chamber of Commerce v.\ Whiting}}, 563 U.S.\ 582 (2011)). Although prolonging a detention solely to verify immigration status would raise serious Fourth Amendment and preemption concerns, state courts could construe \S~2(B) narrowly to require status checks only during an otherwise lawful detention or after release. Under constitutional-avoidance principles (\textit{Fox v.\ Washington}, 236 U.S.\ 273 (1915)), enjoining \S~2(B) prior to state-court interpretation was premature, though as-applied preemption and constitutional challenges remained open.
\end{itemize}

\subsection{Concurring in Part and Dissenting in Part (Justice Scalia)}

Justice Scalia would have upheld all four provisions of S.B.\ 1070. Grounding his opinion in founding-era sovereignty theory (quoting Emer de Vattel, \textit{The Federalist No.\ 42}, the Kentucky and Virginia Resolutions, and \textit{Mayor of New York v.\ Miln}, 11 Pet.\ 102 (1837)), Scalia argued that States possess an inherent sovereign power to exclude unwanted noncitizens unless Congress unequivocally abrogates it or state law directly conflicts with federal law by admitting those federal law excludes or excluding those federal law admits. Implicit field preemption and foreign-affairs sensitivities (citing \textit{Medell\'{i}n v.\ Texas}, 552 U.S.\ 491 (2008)) cannot strip a core attribute of state sovereignty.

Applying that framework, Scalia concurred in upholding \S~2(B) (while faulting the majority for injecting Fourth Amendment concerns into a preemption case) and dissented on the other three sections: \S~6 validly exercises state arrest power and unsolicited ``cooperation'' under \S~1357(g)(10)(B); \S~3 merely enforces federal registration standards rather than adding ``additional or auxiliary'' obligations as in \textit{Hines} (invoking \href{https://supreme.justia.com/cases/federal/us/457/202/}{\textit{Plyler v.\ Doe}}, 457 U.S.\ 202, 228 n.\ 23 (1982), and \textit{California v.\ Zook}, 336 U.S.\ 725 (1949)); and \S~5(C) survives because IRCA's express preemption of employer sanctions (\S~1324a(h)(2)) implies under \textit{expressio unius} that employee sanctions were left to the States, and failed legislative proposals prove nothing. In a pointed closing critique, Scalia cited the Obama Administration's June 15, 2012 Deferred Action for Childhood Arrivals (DACA) memorandum, arguing that executive nonenforcement policies cannot preempt state laws that faithfully enforce the immigration statutes Congress enacted.

\subsection{Concurring in Part and Dissenting in Part (Justice Thomas)}

Justice Thomas likewise would have upheld all four provisions, but rested on a pure textualist rejection of implied ``purposes and objectives'' obstacle preemption rather than Scalia's broader historical theory of state exclusion power. Reaffirming his concurrence in \textit{Wyeth v.\ Levine}, 555 U.S.\ 555, 588 (2009), Thomas contended that the Supremacy Clause gives preemptive force only to the ``ordinary meaning'' of enacted federal statutory text, not to judicially divined congressional purposes. Finding no textual conflict between federal statutes and \S\S~2(B), 3, 5(C), or 6---and rejecting implied field preemption absent express statutory text---Thomas concluded that every challenged provision of S.B.\ 1070 was valid.

\subsection{Concurring in Part and Dissenting in Part (Justice Alito)}

Justice Alito split the difference, agreeing with the majority on \S~2(B) and \S~3 while dissenting on \S~5(C) and \S~6. On \S~2(B), Alito emphasized that Congress invited status inquiries through \S\S~1357(g)(10)(A), 1373(c), and 1644, and called the United States' claim that shifting executive enforcement priorities can preempt state law an ``astounding assertion of federal executive power.'' On \S~3, Alito agreed that \textit{Hines} squarely established field preemption over alien registration.

On \S~5(C), however, Alito faulted the majority for giving ``short shrift'' to the presumption against preemption and \textit{De Canas v.\ Bica}, arguing that employment regulation is a traditional state police power, IRCA's express preemption clause omitted employee sanctions, and Congress's decision not to enact federal criminal penalties does not clearly and manifestly bar States from doing so. On \S~6, Alito argued that the provision added very little to existing Arizona arrest authority, that \S~1357(g)(10)(B) allows unsolicited cooperation, that federal law mandates custody of many criminal aliens under 8 U.S.C.\ \S~1226(c)(1), and that \S~6 should survive a facial challenge under \textit{United States v.\ Salerno}, 481 U.S.\ 739 (1987), because Arizona could confine arrests to \S~1226(c)(1) offenses.

\section{Examples: Future Applications}

\begin{enumerate}
\item \textbf{Same-Side (Preempted): State Criminalization of Unlawful Entry or Reentry.} A border State enacts a statute making unauthorized entry into the State from a foreign country a state misdemeanor and empowering state judges to issue orders requiring noncitizens to return to the foreign nation. \textit{Analysis:} Under \textit{Arizona}, this statute is preempted under both field and obstacle preemption. Congress has enacted a comprehensive federal scheme governing entry and removal (8 U.S.C.\ \S\S~1325, 1326, 1227, 1229a), and unilateral state removal orders directly usurp federal foreign-relations authority and prosecutorial discretion.
\item \textbf{Same-Side (Preempted): Municipal Housing Ordinance Barring Rental to Undocumented Noncitizens.} A city enacts an ordinance requiring residential landlords to verify tenants' immigration status and imposing criminal fines on undocumented individuals who occupy rental housing without federal authorization. \textit{Analysis:} Like S.B.\ 1070's \S\S~3 and 5(C), this ordinance regulates who may remain in the locality based on federal immigration classifications and imposes state-level sanctions to drive out noncitizens whom federal authorities may choose not to remove, obstructing the federal removal framework (consistent with post-\textit{Arizona} appellate rulings such as \textit{Villas at Parkside Partners v.\ City of Farmers Branch}, 726 F.\ 3d 524 (5th Cir.\ 2013)).
\item \textbf{Opposite-Side (Not Preempted): State Prosecution of Identity Theft on Tax Withholding Forms.} A State prosecutes unauthorized workers under a generally applicable state identity-theft statute for using another person's Social Security number on state and federal W-4 tax-withholding forms (without relying on the federal Form I-9). \textit{Analysis:} Unlike \S~5(C), which targeted unauthorized work itself, a neutral state law punishing tax-form identity theft applies to citizens and noncitizens alike and does not conflict with IRCA's employment-verification scheme---the exact distinction the Supreme Court later drew in upholding such prosecutions in \href{https://supreme.justia.com/cases/federal/us/589/17-834/}{\textit{Kansas v.\ Garcia}}, 589 U.S.\ 191 (2020).
\item \textbf{Opposite-Side (Not Preempted): State Business-License Revocation and Mandatory E-Verify for Employers.} A State revokes the business licenses of employers who knowingly hire unauthorized workers and requires all in-state employers to check work eligibility through the federal E-Verify system. \textit{Analysis:} Expressly saved from preemption by IRCA's licensing exception, 8 U.S.C.\ \S~1324a(h)(2), and consistent with Congress's encouragement of E-Verify, as the Court held in \href{https://supreme.justia.com/cases/federal/us/563/582/}{\textit{Chamber of Commerce v.\ Whiting}}, 563 U.S.\ 582 (2011), and reaffirmed in \textit{Arizona}.
\item \textbf{Fence-Sitter (Close Call): State Mandate Requiring Local Jails to Honor 48-Hour ICE Civil Detainers.} A State passes a law prohibiting local sheriffs from adopting ``sanctuary'' policies and requiring county jails to hold any inmate subject to an administrative ICE detainer request for up to 48 hours past the inmate's scheduled state release time so federal officers can assume custody. \textit{Analysis:} On one hand, unlike \S~6's unilateral arrests, honoring an ICE detainer responds directly to a \textit{federal} request, arguably fitting within ``cooperation'' under 8 U.S.C.\ \S~1357(g)(10)(B) and avoiding interference with federal priorities. On the other hand, \textit{Arizona} expressly warned that holding detainees past their authorized release solely for civil immigration purposes without a judicial warrant or formal \S~1357(g) delegation raises severe Fourth Amendment and preemption concerns, leaving the validity of mandatory detainer holds sharply contested in lower courts.
\end{enumerate}

\section{Critique}

\textit{Arizona v.\ United States} drew vigorous criticism from across the ideological spectrum, highlighting enduring tensions in federalism, executive power, and civil rights.

\textbf{Progressive and Civil-Rights Critiques:} Progressive and critical-race scholars criticized the Court's decision to leave Section 2(B) intact and its broader framing of the dispute purely as a federal-state structural contest. In the \href{https://www.stanfordlawreview.org/online/discrimination-preemption-and-arizonas-immigration-law/}{\textit{Stanford Law Review Online}}, Lucas Guttentag argued that by divorcing preemption from antidiscrimination norms (such as 42 U.S.C.\ \S~1981 and \textit{Graham v.\ Richardson}, 403 U.S.\ 365 (1971)), the Court obscured the civil-rights harms of ``attrition through enforcement'' laws. Similarly, commentary in \href{https://www.cambridge.org/core/books/feminist-judgments-immigration-law-opinions-rewritten/commentary-on-arizona-v-united-states-567-us-387-2012/C05E87D040BF471E6BA3D73CFA0BEBE8}{\textit{Feminist Judgments: Immigration Law Opinions Rewritten}} faulted the majority for sidelining the real-world racial profiling directed at Latino residents and citizens. Those concerns bore out in subsequent litigation: as documented in the \href{https://www.law.georgetown.edu/immigration-law-journal/in-print/volume-31-issue-3-spring-2017/arizonas-anti-immigration-law-and-the-pervasiveness-of-racial-profiling/}{\textit{Georgetown Immigration Law Journal}}, civil-rights groups challenged \S~2(B)'s implementation, leading Arizona in 2016 to enter a landmark settlement sharply restricting police officers from prolonging stops or making immigration inquiries based on race or ethnicity.

\textbf{Originalist, Textualist, and Federalism Critiques:} Originalist and textualist scholars echoed Justices Thomas, Scalia, and Alito in attacking the majority's reliance on expansive implied preemption. As summarized in \href{https://www.scotusblog.com/2012/08/academic-highlight-young-on-the-roberts-court-and-preemption/}{commentary on Ernest Young's preemption scholarship} and \href{https://www.scotusblog.com/2011/07/arizona-v-united-states-narrowing-prevent-or-frustrate-conflict-preemption/}{SCOTUSblog's analysis of obstacle preemption}, inferring preemption from congressional silence or failed bills in IRCA undermines the constitutional text of the Supremacy Clause, which makes ``Laws'' enacted through bicameralism and presentment---not uncodified statutory ``purposes'' or shifting executive memoranda---supreme. Critics also noted an internal tension in Kennedy's opinion: while the Court rightly rejected the argument that executive priorities alone could preempt \S~2(B), its invalidation of \S\S~3 and 6 heavily elevated executive prosecutorial discretion into a de facto preemption shield. As noted in the \href{https://www.everycrsreport.com/reports/R42719.html}{Congressional Research Service analysis} and more recent scholarship on \href{https://texaslawreview.org/obsolescent-preemption/}{field preemption in the \textit{Texas Law Review}} and \href{https://scholarship.law.umn.edu/cgi/viewcontent.cgi?article=1189&context=minnlrev}{immigration federalism in the \textit{Minnesota Law Review}}, \textit{Arizona} left unresolved how far States may go when Congress deadlocks on immigration reform and executive administrations swing sharply in enforcement posture---a tension that resurfaced eight years later when the Court narrowed implied preemption in \href{https://supreme.justia.com/cases/federal/us/589/17-834/}{\textit{Kansas v.\ Garcia}}, 589 U.S.\ 191 (2020).

\section{Key Quotations}

All quotations below are verified verbatim from the published opinion in \href{https://supreme.justia.com/cases/federal/us/567/387/}{\textit{Arizona v.\ United States}, 567 U.S.\ 387 (2012)}:

\begin{quote}
``The Government of the United States has broad, undoubted power over the subject of immigration and the status of aliens.'' --- Justice Kennedy, majority opinion, 567 U.S.\ at 394.
\end{quote}

\begin{quote}
``Where Congress occupies an entire field, as it has in the field of alien registration, even complementary state regulation is impermissible. Field preemption reflects a congressional decision to foreclose any state regulation in the area, even if it is parallel to federal standards.'' --- Justice Kennedy, majority opinion, 567 U.S.\ at 401.
\end{quote}

\begin{quote}
``The National Government has significant power to regulate immigration. With power comes responsibility, and the sound exercise of national power over immigration depends on the Nation's meeting its responsibility to base its laws on a political will informed by searching, thoughtful, rational civic discourse. Arizona may have understandable frustrations with the problems caused by illegal immigration while that process continues, but the State may not pursue policies that undermine federal law.'' --- Justice Kennedy, majority opinion, 567 U.S.\ at 416.
\end{quote}

\begin{quote}
``Arizona has moved to protect its sovereignty---not in contradiction of federal law, but in complete compliance with it. The laws under challenge here do not extend or revise federal immigration restrictions, but merely enforce those restrictions more effectively. If securing its territory in this fashion is not within the power of Arizona, we should cease referring to it as a sovereign State.'' --- Justice Scalia, concurring in part and dissenting in part, 567 U.S.\ at 437.
\end{quote}

\begin{quote}
``Under the Supremacy Clause, pre-emptive effect is to be given to congressionally enacted laws, not to judicially divined legislative purposes.'' --- Justice Thomas, concurring in part and dissenting in part, 567 U.S.\ at 440.
\end{quote}

\end{document}
